% ============================================================
%  Analog Color Measurement — Graduate Assignment Writeup
%  Paste into Overleaf (pdfLaTeX, no extra packages needed
%  beyond the standard ones listed below).
% ============================================================
\documentclass[11pt, letterpaper]{article}

% --- Packages -----------------------------------------------
\usepackage[margin=1in]{geometry}
\usepackage{amsmath, amssymb}
\usepackage{graphicx}
\usepackage{booktabs}       % nice tables
\usepackage{hyperref}
\usepackage{caption}
\usepackage{subcaption}
\usepackage{float}
\usepackage{listings}
\usepackage{xcolor}
\usepackage{tikz}
\usetikzlibrary{shapes.geometric, arrows.meta, positioning}

% --- Metadata -----------------------------------------------
\title{\textbf{Analog Color Measurement}\\[4pt]
       \large Capturing and Modeling Hardware Nonlinearities\\
       as Deployable Audio Plugins}
\author{John Micensky}
\date{\today}

% ============================================================
\begin{document}
\maketitle
\tableofcontents
\newpage

% ------------------------------------------------------------
\section{Introduction}
\label{sec:intro}
% TODO: Motivate the problem. Why does analog "color" matter?
% What gap does this tool fill vs. existing impulse-response
% capture workflows (which only model linear behavior)?

% ------------------------------------------------------------
\section{Background and Related Work}
\label{sec:background}
% TODO: Brief survey of:
%   - Hammerstein / Wiener / L-N-L cascade models in audio literature
%   - Existing hardware emulation approaches (convolution, circuit simulation)
%   - Limitations of linear IR capture for saturating devices

% ------------------------------------------------------------
\section{System Overview}
\label{sec:system}

The measurement system is implemented as a JUCE standalone desktop application.
Figure~\ref{fig:system_flow} gives a high-level view of the signal flow from
hardware device under test through to the exported plugin model.

% --- Block diagram placeholder ---
\begin{figure}[H]
    \centering
    % TODO: replace with a real TikZ or inkscape figure
    \fbox{\parbox{0.8\textwidth}{\centering\vspace{1.5cm}
          [System block diagram: Device Under Test $\to$ ADC $\to$
           Capture Pipeline $\to$ Analysis Engine $\to$ LNL Model $\to$ Plugin]\vspace{1.5cm}}}
    \caption{End-to-end signal flow of the measurement and modeling pipeline.}
    \label{fig:system_flow}
\end{figure}

\subsection{Application Architecture}
% TODO: Describe the major software components:
%   MainComponent (UI), AudioEngine (real-time I/O),
%   StimulusPlan (test sequence), SessionWriter (disk I/O),
%   AnalysisEngine (offline), LNLModel (data structure), Plugin target.

\subsection{Session Outputs}
Each measurement session produces four artifacts per stimulus step:
\begin{itemize}
    \item \texttt{*\_ref.wav} — the direct reference signal recorded simultaneously with playback
    \item \texttt{*\_rec.wav} — the return signal captured from the hardware device
    \item \texttt{session.json} — metadata (device name, sample rate, gain labels, timestamps)
    \item \texttt{markers.json} — per-step index mapping filenames to gain settings and stimuli
\end{itemize}

% ------------------------------------------------------------
\section{Stimulus Design}
\label{sec:stimulus}

Measurements are organized into a \emph{stimulus plan}: an ordered sequence of
$(g, s)$ pairs where $g$ is a hardware gain setting and $s$ is one of the
available test stimuli.  Three quality tiers are defined:

\begin{table}[H]
    \centering
    \begin{tabular}{lll}
        \toprule
        \textbf{Quality} & \textbf{Stimuli} & \textbf{Purpose} \\
        \midrule
        Quick    & Sine sweep, 1\,kHz sine                                      & Fast sanity check \\
        Standard & Sine sweep, 1\,kHz sine, 100\,Hz sine, drums excerpt         & Balanced coverage \\
        HiFi     & Sine sweep, pink noise, 1\,kHz, 100\,Hz, drums excerpt       & Full characterization \\
        \bottomrule
    \end{tabular}
    \caption{Stimulus quality tiers and their constituent signals.}
    \label{tab:quality}
\end{table}

% TODO: Add rationale for each stimulus choice (e.g., why 100 Hz specifically,
% why a real drums excerpt rather than a second sine).

% ------------------------------------------------------------
\section{Latency Alignment}
\label{sec:latency}

Before any analysis, the recorded signal $r[n]$ must be time-aligned with the
reference signal $x[n]$ to compensate for round-trip hardware latency.
The integer-sample delay $\hat{\tau}$ is estimated via normalized
cross-correlation:

\begin{equation}
    \hat{\tau} \;=\; \arg\max_{\tau} \;
    \frac{\displaystyle\sum_{n} x[n]\, r[n - \tau]}
         {\sqrt{\displaystyle\sum_{n} x[n]^{2}\;
                \displaystyle\sum_{n} r[n]^{2}}}
    \label{eq:xcorr}
\end{equation}

The alignment is computed efficiently in the frequency domain using the FFT:

\begin{equation}
    C[k] \;=\; X[k] \cdot R^{*}[k],
    \qquad
    c[n] \;=\; \mathcal{F}^{-1}\!\left\{C[k]\right\},
    \qquad
    \hat{\tau} \;=\; \arg\max_{n}\; c[n]
    \label{eq:xcorr_fft}
\end{equation}

% TODO: Report measured latency values on your test interface.
% TODO: Discuss sub-sample interpolation if implemented.

% ------------------------------------------------------------
\section{Gray-Box L--N--L Model}
\label{sec:lnl}

\subsection{Model Motivation}
\label{sec:lnl_motivation}

A purely \emph{black-box} model (e.g., a deep neural network trained end-to-end)
can achieve high accuracy but offers no interpretability and is expensive to
deploy in real-time audio contexts.  A purely \emph{white-box} model (e.g., a
full circuit simulation) is physically accurate but requires proprietary
schematic knowledge and is computationally prohibitive at audio rates.

This project adopts a \emph{gray-box} approach: the model topology is fixed by
prior knowledge of analog signal chains (a cascade of linear filtering and a
static nonlinearity), while the parameters of each block are estimated
empirically from the captured data.

\subsection{Model Structure}
\label{sec:lnl_structure}

The device under test is modeled as a \emph{Linear--Nonlinear--Linear} (L--N--L)
cascade, illustrated in Figure~\ref{fig:lnl_diagram}.

\begin{figure}[H]
\centering
\begin{tikzpicture}[
    block/.style = {rectangle, draw, minimum width=2.2cm, minimum height=1cm,
                    text centered, font=\small},
    line/.style  = {draw, -{Stealth[length=2mm]}}
]
    \node[block]                       (l1) {$L_1[z]$\\{\tiny Input Filter}};
    \node[block, right=1.4cm of l1]    (n)  {$\mathcal{N}(\cdot)$\\{\tiny Waveshaper}};
    \node[block, right=1.4cm of n]     (l2) {$L_2[z]$\\{\tiny Output Filter}};

    \coordinate (in)  at ([xshift=-1.4cm]l1.west);
    \coordinate (out) at ([xshift=+1.4cm]l2.east);

    \line (in)    -- node[above]{\small$x[n]$} (l1.west);
    \line (l1.east) -- node[above]{\small$v[n]$} (n.west);
    \line (n.east)  -- node[above]{\small$w[n]$} (l2.west);
    \line (l2.east) -- node[above]{\small$\hat{y}[n]$} (out);
\end{tikzpicture}
\caption{L--N--L gray-box model. The input signal $x[n]$ passes through a
linear FIR filter $L_1$, a static memoryless nonlinearity $\mathcal{N}$,
and a second linear FIR filter $L_2$ to produce the model output
$\hat{y}[n]$.}
\label{fig:lnl_diagram}
\end{figure}

The model output is:

\begin{equation}
    \hat{y}[n] \;=\; \bigl(h_2 * \mathcal{N}(h_1 * x)\bigr)[n]
    \label{eq:lnl}
\end{equation}

where $h_1$ and $h_2$ are the FIR impulse responses of $L_1$ and $L_2$
respectively, $*$ denotes linear convolution, and $\mathcal{N}$ is a
point-wise static function.

\subsection{Block Identification}
\label{sec:lnl_id}

Each block is identified independently from dedicated measurements.

\paragraph{$L_1$: Input Filter.}
The linear frequency response of the device is measured using a sine sweep
stimulus.  The transfer function magnitude is estimated bin-wise as:

\begin{equation}
    \bigl|\hat{H}(f_k)\bigr| \;=\;
    \frac{\bigl|\mathcal{F}\{r\}[k]\bigr|}{\bigl|\mathcal{F}\{x\}[k]\bigr|}
    \label{eq:fr_estimate}
\end{equation}

where $x$ and $r$ are the (aligned) reference and recorded signals.
The magnitude spectrum is 1/3-octave smoothed and used to design a
$512$-tap linear-phase FIR filter $h_1$ via an inverse-FFT window method.

\paragraph{$\mathcal{N}$: Static Nonlinearity (Waveshaper).}
With $L_1$ applied, the pre-filtered signal $v[n] = (h_1 * x)[n]$ is
compared to the recorded output at each sample amplitude.  The mapping
$v[n] \mapsto r[n]$ is accumulated into a 1024-entry lookup table
(waveshaper) uniformly spanning the input range $[-1,\,+1]$:

\begin{equation}
    \mathcal{N}(a) \;\approx\; \text{LUT}\!\left[
        \left\lfloor \frac{(a + 1)}{2} \cdot 1023 \right\rceil
    \right]
    \label{eq:waveshaper}
\end{equation}

Measurements at multiple gain levels and stimulus frequencies (100\,Hz and
1\,kHz sine tones) are pooled to improve coverage of the input amplitude range.

\paragraph{$L_2$: Output Filter.}
In the current implementation $L_2$ is the identity ($h_2 = \delta$).
A non-trivial output filter may be fitted in future work to capture
output-side coloration that is not attributable to the input chain or
the nonlinearity.

\subsection{Total Harmonic Distortion}
\label{sec:thd}

As a diagnostic, Total Harmonic Distortion (THD) is computed for each
sine-tone capture:

\begin{equation}
    \text{THD} \;=\;
    \frac{\sqrt{\displaystyle\sum_{k=2}^{K} A_k^2}}{A_1} \times 100\%
    \label{eq:thd}
\end{equation}

where $A_1$ is the fundamental amplitude and $A_k$ is the amplitude of the
$k$-th harmonic, all estimated from the magnitude spectrum of the recorded signal.
THD is reported per gain label to characterize the onset and degree of saturation.

% ------------------------------------------------------------
\section{Measurement Procedure}
\label{sec:procedure}
% TODO: Step-by-step instructions a reader would follow to
% characterize a new device. Reference the UI workflow.

% ------------------------------------------------------------
\section{Plugin Implementation}
\label{sec:plugin}
% TODO: Describe how the fitted LNLModel is exported to a JUCE
% AudioProcessor. Discuss real-time convolution strategy for L1
% (overlap-add/save) and LUT interpolation for N.

% ------------------------------------------------------------
\section{Results}
\label{sec:results}
% TODO: Apply to a specific hardware device.
% Include FR plot, THD vs. gain table, waveshaper curve,
% and A/B perceptual evaluation.

\subsection{Frequency Response}
% TODO: Figure — measured vs. modeled FR.

\subsection{Nonlinearity Characterization}
% TODO: Figure — waveshaper curve + THD table.

\subsection{Perceptual Evaluation}
% TODO: Informal listening test methodology and findings.

% ------------------------------------------------------------
\section{Discussion}
\label{sec:discussion}
% TODO: Model limitations (memory effects, dynamic nonlinearities,
% compressor behavior not yet captured). Future work.

% ------------------------------------------------------------
\section{Conclusion}
\label{sec:conclusion}
% TODO: Summarize contributions and outcomes.

% ------------------------------------------------------------
\bibliographystyle{plain}
\bibliography{refs}
% TODO: populate refs.bib with relevant citations.

\end{document}
